%% summaries.tex -- pushdown summaries of the covariant signed system: the
%% lowest point, duality and sign coherence, least closures of regular seeds,
%% forbidden roots and the pop bound, deadness (by support, or by a self event),
%% closure membership for a live seed, head prefixes, and the decision of
%% satisfiability and of (I0). Long proofs: appendix-summaries.tex (app:sm).

\section{Pushdown Summaries}\label{sec:sm}

The three conditions of \cref{thm:ch-characterization} concern infinite sets
of atoms: least closures, the sets $\Reach_\chi(a)$, and live atoms. This section
describes them with finite automata over the alphabet of directions. The pure
moves of a system form a pushdown system whose control states are the signed
variables and whose stack is the word of an atom (\cref{sec:lf-clauses}). They
are the moves of the signed system $\chi^\pm$, a covariant system, so the
upper and lower bound automata of covariant systems summarize the parts of a
pure derivation that pop and push its stack, and every pure derivation splits
at its lowest point into one part of each kind (\cref{sec:sm-lowest}). With
them we describe least closures, also of regular sets of seeds
(\cref{sec:sm-regular}).

Two consequences of the signs have no counterpart for a covariant
constructor. First, the lower automata are the upper automata with all signs
flipped (\cref{lem:sm-duality}). Second, an atom can be dead because the
closure of its seed contains the opposite of an extension of the seed, the
self-collision of \cref{cor:lf-dead}(c). We show that a self-collision of an atom
that is not already dead by support is found by comparing two suffixes of the
seed word: it yields a strict \emph{self event}, and deadness is the union of
deadness by support, which is regular, and of the self events
(\cref{thm:sm-dead}). Deadness need not be regular (\cref{ex:sm-arrow}). The
section ends with closure membership for a live seed (\cref{lem:sm-membership}), the heads of pure derivations
(\cref{lem:sm-head-prefix}), and the decision of satisfiability and of
condition \CondI\ (\cref{prop:sm-decide}). Proofs omitted here are in
\cref{app:sm}.

\paragraph{Conventions}
Throughout, $\chi$ is a system over a finite set $V$ of variables, we use the
notation of \cref{sec:bg,sec:lf} and write $\bar P=\{\bar p:p\in P\}$ for
$P\subseteq V^\pm$, and for $n=0$, where $\eps$ is the only word, every
statement below holds trivially.

\subsection{Summaries and the Lowest Point}\label{sec:sm-lowest}

By \cref{sec:lf-clauses}, the pure moves of $\chi$ are the moves of the copy
and constructor literals of $\chi^\pm$, read as a covariant system over
$V^\pm$: a literal $p\le q$ gives $(p,w)\to(q,w)$, a literal $f(p_i)_i\le q$
gives $(p_i,w)\to(q,iw)$, and a literal $q\le f(p_i)_i$ gives
$(q,iw)\to(p_i,w)$; and the forbidden roots are the atoms $(q,\eps)$ for the
roots $q$ of the upper constructor literals of $\chi^\pm$. The extreme
literals of $\chi^\pm$ give no moves. The \emph{height} of an atom $(p,w)$ is
$\len w$, and a pure derivation is \emph{above} a word $\sigma$ if every atom
on it has the form $(p,\eta\sigma)$. The automata below are the bound automata
of a covariant system~\citep{Zhao26}, for an arbitrary
arity;\footnote{In the notation of \cref{app:bg-sat}, $D_\chi$ is the closure
$D^{\mathrm{cov}}_{\chi^\pm}$, with the copy pairs included in the
saturation.} we write $\chi$ instead of $\chi^\pm$ in subscripts.

The closure $D_\chi$ summarizes the pure derivations that return to their
starting height without going below it. Besides copies, such a derivation from
$a_i$ to $b_i$ can push the letter $i$ through a literal $f(a_j)_j\le u$,
continue with a derivation of the same kind from $u$ to $v$, one level higher,
and pop $i$ again through a literal $v\le f(b_j)_j$; this is the decomposition
rule below, and transitivity strings such pieces together. The upper automaton of $p$, after reading $\gamma$, holds the signed
variables reached from $(p,\gamma\sigma)$ by popping $\gamma$, and the lower
automaton of $q$, after reading $\beta$, those from which $(q,\beta\sigma)$ is
reached by pushing $\beta$ (\cref{lem:sm-summaries}).

\begin{definition}[Closure and bound automata]\label{def:sm-automata}
The \emph{closure} $D_\chi$ is the least reflexive and transitive relation on
$V^\pm$ that contains every pair $(p,q)$ with $p\le q$ in $\chi^\pm$ and
satisfies the decomposition rule
\[
 \frac{f(a_j)_j\le u\in\chi^\pm\qquad D_\chi(u,v)\qquad v\le f(b_j)_j\in\chi^\pm}
      {D_\chi(a_i,b_i)}\qquad(i\in\Sigma).
\]
Put $\upc\chi p=\{r:D_\chi(p,r)\}$ and $\downc\chi q=\{r:D_\chi(r,q)\}$. For
$P,Q\subseteq V^\pm$ and $i\in\Sigma$, let
\[
\begin{aligned}
 \pstep i(P)&=\{b:\ \exists r\in P,\ \exists\,(r\le f(b_j)_j)\in\chi^\pm,\
  D_\chi(b_i,b)\},\\
 \qstep i(Q)&=\{a:\ \exists r\in Q,\ \exists\,(f(a_j)_j\le r)\in\chi^\pm,\
  D_\chi(a,a_i)\}.
\end{aligned}
\]
The \emph{upper automaton} of $p\in V^\pm$ has state set $2^{V^\pm}$, start
state $\upc\chi p$, and transitions $\pstep i$; its state after a word $w$ is
$P^\chi_p(w)$. The \emph{lower automaton} of $q$ has start state $\downc\chi
q$ and transitions $\qstep i$; its state after $w$ is $Q^\chi_q(w)$. At cuts we
write $P^\chi_{p,j}(w)=P^\chi_p(\wpre wj)$ and $Q^\chi_{q,j}(w)=Q^\chi_q(\wpre wj)$.
\end{definition}

Both automata are deterministic and complete, and they read their word from
left to right: $P^\chi_p(wi)=\pstep i(P^\chi_p(w))$. They have $4^{\len V}$
states each, and $D_\chi$ and every transition are computable in time
polynomial in the size of $\chi$. Every set $P^\chi_p(w)$ is upward closed
under $D_\chi$ and every $Q^\chi_q(w)$ is downward closed, the empty set is
absorbing in both automata, and $\pstep i$ and $\qstep i$ distribute over
unions.

\begin{lemma}[Pushdown summaries]\label{lem:sm-summaries}
For all signed variables $a,b,p,q$ and words $w,\beta,\gamma,\sigma$:
\begin{enumerate}[label=(\alph*),leftmargin=*]
\item $D_\chi(a,b)$ iff $(a,w)\to^*_\chi(b,w)$ by a pure derivation above
  $w$;
\item $b\in P^\chi_p(\gamma)$ iff $(p,\gamma\sigma)\to^*_\chi(b,\sigma)$ by a
  pure derivation above $\sigma$;
\item $a\in Q^\chi_q(\beta)$ iff $(a,\sigma)\to^*_\chi(q,\beta\sigma)$ by a
  pure derivation above $\sigma$.
\end{enumerate}
In particular, the existence of these derivations does not depend on $w$ or
$\sigma$.
\end{lemma}

\begin{figure}[tbp]
\centering
\begin{tikzpicture}[x=1.15cm,y=1.0cm,>=Stealth,font=\small]
  \fill[black!8] (0,0) rectangle (6.4,1);
  \node[font=\footnotesize,anchor=west] at (0.1,0.5)
    {suffix $\sigma$ of every word, never read};
  \draw[->] (0,0) -- (6.6,0) node[below left,font=\footnotesize] {moves};
  \draw[->] (0,0) -- (0,3.5) node[left,font=\footnotesize] {height};
  \draw[dashed] (0,1) -- (6.4,1);
  \node[left,font=\footnotesize] at (0,1) {$\len\sigma$};
  \node[left,font=\footnotesize] at (0,3) {$\len w$};
  \node[left,font=\footnotesize] at (0,2.5) {$\len v$};
  \draw[thick] (0,3) -- (0.4,2.4) -- (0.8,2.7) -- (1.3,2.0) -- (1.8,1.6)
    -- (2.2,1.9) -- (2.7,1.3) -- (3.1,1) -- (3.5,1.4) -- (3.9,1.2)
    -- (4.4,1.7) -- (4.9,2.1) -- (5.3,1.9) -- (5.8,2.5);
  \fill (0,3) circle (1.5pt) node[above right] {$(p,\gamma\sigma)$};
  \fill (3.1,1) circle (1.5pt);
  \node[above,font=\footnotesize] at (3.1,1.05) {$(r,\sigma)$};
  \fill (5.8,2.5) circle (1.5pt) node[above] {$(q,\beta\sigma)$};
  \draw[decorate,decoration={brace,amplitude=4pt,mirror}] (0,-0.35) -- (3.1,-0.35)
    node[midway,below=4pt,font=\footnotesize,align=center]
    {pops $\gamma$\\ $r\in P^\chi_{p}(\gamma)$};
  \draw[decorate,decoration={brace,amplitude=4pt,mirror}] (3.1,-0.35) -- (5.8,-0.35)
    node[midway,below=4pt,font=\footnotesize,align=center]
    {pushes $\beta$\\ $r\in Q^\chi_{q}(\beta)$};
\end{tikzpicture}
\caption{The lowest point of a pure derivation from $(p,w)$ to $(q,v)$
(\cref{lem:sm-lowest-point}). No atom on the derivation goes below the least
height $\len\sigma$, so the suffix $\sigma$ is never read: $w=\gamma\sigma$ and
$v=\beta\sigma$. The first atom $(r,\sigma)$ of least height splits the
derivation into a part summarized by the upper automaton of $p$ and a part
summarized by the lower automaton of $q$.}
\Description{A graph of word length against moves. The curve starts at the
height of w, wanders down to a lowest point at the height of sigma, and climbs
back to the height of v. A shaded band below the lowest height marks the suffix
sigma, which is never read.}
\label{fig:sm-lowest-point}
\end{figure}

\begin{lemma}[Lowest point]\label{lem:sm-lowest-point}
For signed variables $p,q$ and words $w,v$, $(p,w)\to^*_\chi(q,v)$ iff there
are cuts $i\le\len w$ and $e\le\len v$ with $\wsuf wi=\wsuf ve$ and
$P^\chi_{p,i}(w)\cap Q^\chi_{q,e}(v)\ne\varnothing$.
\end{lemma}

Cut the derivation at its first atom of least height
(\cref{fig:sm-lowest-point}) and apply \cref{lem:sm-summaries}(b,c) to the two
parts; the signed variable $r$ at the lowest point may belong to a visible or
to a hidden variable, with either sign. In the notation of
\cref{fig:sm-lowest-point}, $\gamma$ is the \emph{pop head} and $\beta$ the
\emph{push head} of the derivation.

\begin{lemma}[Duality and sign coherence]\label{lem:sm-duality}
For all signed variables $a,b,p,q$ and words $w$:
\begin{enumerate}[leftmargin=*]
\item $D_\chi(a,b)$ iff $D_\chi(\bar b,\bar a)$; hence
  $\downc\chi{\bar q}=\overline{\upc\chi q}$;
\item $Q^\chi_{\bar q}(w)=\overline{P^\chi_q(w)}$;
\item $D_\chi(a,b)$ implies $\sgn(a)=\sgn(b)$; every $r\in P^\chi_p(w)$ has
  $\sgn(r)=\sgn(p)\oplus\parity(w)$; and every $r\in Q^\chi_q(w)$ has
  $\sgn(r)=\sgn(q)\oplus\parity(w)$.
\end{enumerate}
\end{lemma}

Items (1) and (2) reverse pure derivations by duality (\cref{lem:lf-duality}):
the reversed derivation visits the same words, so it stays above the same
suffix. By item (2), \cref{lem:sm-lowest-point} can be read with upper
automata only. Item (3) is \emph{sign coherence}, the sign invariant of
\cref{lem:lf-moves}(2) applied to the derivations of \cref{lem:sm-summaries}.

\begin{example}[Shift closure: the signed system]\label{ex:sm-running}
Let $n=2$ with $f$ covariant, and let
$\psi=\{x\le z,\ z\le y,\ z\le f(z,z)\}$ be the consequent of
\cref{ex:intro-shift}, now over $\BFT$, with $z$ hidden. The copy and
constructor literals of its signed system are
\[
 x^0\le z^0,\quad z^0\le y^0,\quad z^0\le f(z^0,z^0),\quad
 y^1\le z^1,\quad z^1\le x^1,\quad f(z^1,z^1)\le z^1 .
\]
It has no facts, and its only forbidden root is $(z^0,\eps)$: the
root of $z$ is not $\Top$. Since $f$ is covariant, $(x^0,w)$ says that $\rho x$
has a leaf $\Top$ at a prefix of $w$, and $(x^1,w)$ that it has a leaf $\Bot$
at a prefix of $w$. The decomposition rule never applies, so $D_\psi$ is the
reflexive and transitive closure of the copies. Hence
$\upc\psi x^0=\{x^0,z^0,y^0\}$ and $P^\psi_{x^0}(w)=\{z^0,y^0\}$ for
$w\ne\eps$; dually, $\downc\psi x^1=\{x^1,z^1,y^1\}$ and
$Q^\psi_{x^1}(w)=\{z^1,y^1\}$ for $w\ne\eps$, as \cref{lem:sm-duality}(2)
predicts. By \cref{lem:sm-lowest-point}, $(y^1,\sigma)\to^*_\psi(x^1,\beta\sigma)$
for all words $\beta$ and $\sigma$: in every solution of $\psi$, if $y$ has a
leaf $\Bot$ at a prefix of $\sigma$, then $x$ has a leaf $\Bot$ at a prefix of
$\beta\sigma$. For trees over $\BF$ this says that every suffix of a
constructor position of $x$ is a constructor position of $y$. Similarly
$(x^0,w)\to^*_\psi(z^0,\eps)$ for every $w$, a pure derivation to the forbidden
root: $\psi$ forbids leaves $\Top$ in $x$.
\end{example}

\subsection{Least Closures of Regular Seeds}\label{sec:sm-regular}

Let
\[
 \Fup\chi=\{r:\ r\in P^\chi_e(\gamma)\text{ for some }(e,\eps)\in\Fact_\chi
 \text{ and some word }\gamma\}
\]
be the set of signed variables that a fact reaches at the empty word
after popping a word.

\begin{lemma}[The least closure]\label{lem:sm-mu}
For every signed variable $q$ and word $v$,
\[
 (q,v)\in\mu_\chi\iff Q^\chi_{q,e}(v)\cap\Fup\chi\ne\varnothing
 \text{ for some cut }e\le\len v .
\]
The set $\Fup\chi$ is the least set that contains $\upc\chi p$ for every
fact $(p,\eps)$ and contains $\pstep i(\{r\})$ for each of its
elements $r$ and each $i\in\Sigma$; it is computable in polynomial time.
\end{lemma}

A fact $(p,\eps)$ derives $(r,\sigma)$ for every $\sigma$: extend it
to $(p,\gamma\sigma)$ and pop $\gamma$. So $\mu_\chi$ contains every atom that a
pure derivation pushing a word reaches from such an $(r,\sigma)$, and
\cref{lem:lf-ext-first,lem:sm-lowest-point} show that there is nothing else.
The same argument applies to any regular set of seeds.

\begin{lemma}[Closure of regular seeds]\label{lem:sm-regular-closure}\label{prop:sm-regular-closure}
Let $S$ be a set of atoms whose slices $S_p=\{w:(p,w)\in S\}$ are regular for
all $p\in V^\pm$. For signed variables $p,r,q$, put
$U_{p,r}=\{\gamma:r\in P^\chi_p(\gamma)\}$,
$V_{r,q}=\{\beta:r\in Q^\chi_q(\beta)\}$, and
$\mathrm{Res}_{p,r}=\{\sigma:\gamma\sigma\in S_p\Sigma^*\text{ for some }
\gamma\in U_{p,r}\}$. Then the slice of $H_\chi(S)$ at $q$ is
\[
 H_\chi(S)_q=\bigcup_{r\in\Fup\chi}V_{r,q}\,\Sigma^*\ \cup\
 \bigcup_{p,r\in V^\pm}V_{r,q}\,\mathrm{Res}_{p,r}.
\]
All these languages are regular, and finite automata for them are computable
from finite automata for the slices of $S$. Consequently, whether $H_\chi(S)$
is consistent is decidable.
\end{lemma}

In particular, $\mu_\chi$ (the case $S=\varnothing$) and the closures of one
or two seeds have regular slices, and their consistency is decidable. The
lemma does not say that the set of live atoms is regular, which is
false in general (\cref{ex:sm-arrow}).

\subsection{Forbidden Roots and the Pop Bound}\label{sec:sm-cone}

A pure derivation pops a letter only at the root of an upper constructor
literal of $\chi^\pm$, and every such root is forbidden at the empty word. This
gives the \emph{pop bound}, the counterpart of the cone lemma of a covariant
system~\citep{Zhao26}, stated with the following notion and with
$\Ndown\chi=\{r:D_\chi(r,m)\text{ for some }(m,\eps)\in\Neg_\chi\}$.

\begin{definition}[Root deadness]\label{def:sm-deadu}
An atom $(p,w)$ is \emph{root-dead} in $\chi$, written $\DeadU_\chi(p,w)$, if
$\Reach_\chi(p,w)$ contains a forbidden root of $\chi$.
\end{definition}

\begin{lemma}[Roots and the pop bound]\label{lem:sm-roots}
For every signed variable $p$ and all words $w,v$:
\begin{enumerate}[leftmargin=*]
\item if $P^\chi_p(wi)\ne\varnothing$ for some letter $i$, then some
  $r\in P^\chi_p(w)$ is the root of an upper constructor literal of
  $\chi^\pm$, $(p,w)\to^*_\chi(r,\eps)$, and $(r,\eps)\in\Neg_\chi$; in
  particular $\DeadU_\chi(p,w)$;
\item \emph{(pop bound)} if $\DeadU_\chi(p,\wpre wk)$ fails, then
  $P^\chi_{p,j}(w)=\varnothing$ for every cut $j$ with $k<j\le\len w$;
\item $\DeadU_\chi(p,w)$ iff $P^\chi_p(wz)\cap\Ndown\chi\ne\varnothing$ for
  some word $z$;
\item if $\DeadU_\chi(p,wv)$, then $\DeadU_\chi(p,w)$.
\end{enumerate}
\end{lemma}

\begin{proof}
(1) If $b\in\pstep i(P^\chi_p(w))$, some $r\in P^\chi_p(w)$ is the root of an
upper constructor literal $r\le f(b_j)_j$ of $\chi^\pm$; then
$(r,\eps)\in\Neg_\chi$ by \cref{sec:lf-clauses}, and
$(p,w)\to^*_\chi(r,\eps)$ by \cref{lem:sm-summaries}(b) with $\sigma=\eps$.
(2) If $P^\chi_{p,j}(w)\ne\varnothing$ with $j>k$, then
$P^\chi_{p,k+1}(w)\ne\varnothing$, because the empty set is absorbing, and (1)
applied to $\wpre wk$ and the letter $w[k]$ gives $\DeadU_\chi(p,\wpre wk)$.
(3) By \cref{lem:lf-ext-first}, $(m,\eps)\in \Reach_\chi(p,w)$ iff
$(p,w\alpha)\to^*_\chi(m,\eps)$ for some $\alpha$; by
\cref{lem:sm-lowest-point} with the empty word $v=\eps$, whose only cut is $0$,
this holds iff $P^\chi_p(w\alpha)\cap\downc\chi m\ne\varnothing$. (4) By
\cref{lem:lf-closure}(4), $(p,wv)\in \Reach_\chi(p,w)$, so
$\Reach_\chi(p,wv)\subseteq \Reach_\chi(p,w)$.
\end{proof}

\subsection{Deadness}\label{sec:sm-dead}

\Cref{cor:lf-dead} lists three causes of deadness for a satisfiable system:
(a) a forbidden root in $\Reach_\chi(p,w)$, which is $\DeadU_\chi(p,w)$; (b) the
opposite of an extension of the seed in $\mu_\chi$; and (c) a self-collision,
the opposite of an extension of the seed in $\Reach_\chi(p,w)$. The causes may
overlap. Causes (a) and (b) together are deadness by support, which automata
read. A self-collision of an atom that is not dead by support is read by a
strict comparison of the seed word.

\begin{definition}[Deadness by support, self events]\label{def:sm-dead}
An atom $(p,w)$ is \emph{dead by support} in $\chi$, written
$\DeadT_\chi(p,w)$, if $\DeadU_\chi(p,w)$ or $(\bar p,w\delta)\in\mu_\chi$ for
some word $\delta$. A \emph{self event} of $(p,w)$ in $\chi$ is a comparison
$(i,j)$ of $w$ (\cref{def:bg-comparison}), that is, a pair of cuts
$0\le i<j\le\len w$ with $\wsuf wj\pre\wsuf wi$, such that
\[
 P^\chi_{p,i}(w)\cap\overline{P^\chi_{p,j}(w)}\ne\varnothing .
\]
Self events are strict: $i<j$. We write $\SelfEv_\chi(p,w,i,j)$ if $(i,j)$ is
a self event of $(p,w)$.
\end{definition}

Self events have a counterpart in the plain-entailment proof, whose readers
already compare a root with its own signed copy \citep{Zhao26}. Here
both cuts belong to the word of one seed: the state of the upper automaton of
$p$ at the cut $i$ and the flipped state at the cut $j$ share a signed
variable, which serves as the lowest point of a pure derivation from an
extension of $(p,w)$ to the opposite of an extension of $(p,w)$
(\cref{thm:sm-dead}).

Both kinds of deadness by support are read by the upper automaton of $p$
together with one bit. Put
\[
 \DeadU^\eps_\chi=\{r\in V^\pm:\DeadU_\chi(r,\eps)\},\qquad
 \DeadT^\eps_\chi=\{r\in V^\pm:\DeadT_\chi(r,\eps)\}.
\]

\begin{lemma}[Deadness from states]\label{lem:sm-dead-states}
For every signed variable $p$ and all words $w$, $v$:
\begin{enumerate}[leftmargin=*]
\item $\DeadU_\chi(p,w)$ iff $P^\chi_p(w)\cap\DeadU^\eps_\chi\ne\varnothing$;
\item $(\bar p,w)\in\mu_\chi$ iff
  $\overline{P^\chi_{p,e}(w)}\cap\Fup\chi\ne\varnothing$ for some cut
  $e\le\len w$;
\item $\DeadT_\chi(p,w)$ iff $(\bar p,w)\in\mu_\chi$ or
  $P^\chi_p(w)\cap\DeadT^\eps_\chi\ne\varnothing$;
\item if $\DeadT_\chi(p,wv)$, then $\DeadT_\chi(p,w)$;
\item the sets $\DeadU^\eps_\chi$ and $\DeadT^\eps_\chi$ are computable in
  polynomial time.
\end{enumerate}
\end{lemma}

So the words $w$ with $\DeadU_\chi(p,w)$, and those with $\DeadT_\chi(p,w)$,
form regular languages that are closed under prefixes: the first is read by
the upper automaton of $p$, and the second by that automaton together with a
bit that records whether some state $P$ read so far has
$\bar P\cap\Fup\chi\ne\varnothing$. Self events, in contrast, compare two cuts
of one word. They are what remains of deadness.

\begin{theorem}[Deadness]\label{thm:sm-dead}
Let $\chi$ be satisfiable. An atom $(p,w)$ is dead iff $\DeadT_\chi(p,w)$ or
$(p,w)$ has a self event in $\chi$. More precisely, every self event of $(p,w)$
yields a self-collision of $(p,w)$, and if $(p,w)$ has a self-collision but
$\DeadT_\chi(p,w)$ fails, then $(p,w)$ has a self event.
\end{theorem}

By \cref{cor:lf-dead}, it remains to treat a self-collision, a pure
derivation from an extension $(p,w\alpha)$ to an opposite atom
$(\bar p,w\delta)$, of an atom $(p,w)$ that is not dead by support. Split it
at its lowest point. If its pop head or its push head is longer than $w$, then the upper automaton of $p$ is
nonempty after some proper extension of $w$ (for the push head, by
\cref{lem:sm-duality}(2)), and the pop bound (\cref{lem:sm-roots}) gives
$\DeadU_\chi(p,w)$. Otherwise both heads end at cuts $i,j$ of $w$, the
suffixes $\wsuf wi$ and $\wsuf wj$ are comparable in the prefix order, and by
\cref{lem:sm-duality}(2) the lowest point lies in
$P^\chi_{p,i}(w)\cap\overline{P^\chi_{p,j}(w)}$. Sign coherence excludes
$i=j$, and the condition is symmetric in $i$ and $j$, so the shorter suffix is
a prefix of the longer one and $(i,j)$ or $(j,i)$ is a strict self event.
Conversely, every self event yields a self-collision. The proof is in
\cref{app:sm}.

\begin{example}[A nonregular dead language]\label{ex:sm-arrow}
Let $\varphi=\{x\le x\to x\}$ for the arrow, as in
\cref{ex:lf-clauses,ex:lf-least}. The copy and constructor literals of
$\varphi^\pm$ are $x^0\le f(x^1,x^0)$ and $f(x^0,x^1)\le x^1$, there are no
facts, so $\mu_\varphi=\varnothing$, and the only forbidden root is
$(x^0,\eps)$. The closure $D_\varphi$ is the identity, and
\[
 P^\varphi_{x^0}(1^k)=\{x^0\},\qquad P^\varphi_{x^0}(1^k0)=\{x^1\},\qquad
 P^\varphi_{x^0}(w)=\varnothing\ \text{if $1^k0$ is a proper prefix of $w$}.
\]
Hence $\DeadT^\eps_\varphi=\{x^0\}$, and $\DeadT_\varphi(x^0,w)$ holds iff
$w\in1^*$. The guard of a self event $(i,j)$ requires
$P^\varphi_{x^0,i}(w)=\{x^0\}$ and $P^\varphi_{x^0,j}(w)=\{x^1\}$, so
$w=1^k0v$, $i\le k$, and $j=k+1$; the comparison requires $v\pre1^m0v$ with
$m=k-i$, that is, $v\pre(1^m0)^\omega$ (\cref{lem:bg-periodic}). By
\cref{thm:sm-dead}, the dead words of $x^0$ are
\[
 1^*\ \cup\ \{1^k0v:\ v\pre(1^m0)^\omega\text{ for some }m\le k\}.
\]
This language is not regular: its intersection with $1^*01^*0$ is
$\{1^k01^m0:m\le k\}$. For example, $(x^0,0)$ is dead by the self event
$(0,1)$ of the word $0$, although $\DeadT_\varphi(x^0,0)$ fails. This is the
self event that the proof of \cref{thm:sm-dead} extracts from the
self-collision $(x^0,00)\to(x^1,0)$ of \cref{ex:lf-least}: its lowest point
$(x^1,0)$ lies at the cut $1$ of $00$ and the cut $0$ of $0$. The atom
$(x^0,01)$ is live (\cref{ex:lf-least}).
\end{example}

\begin{remark}[No self events without mixed parity]\label{rem:sm-parity}
If $\chi$ has a parity assignment $\pi$ (\cref{rem:ch-parity}), then every
pure move preserves $\sgn(r)\oplus\pi(r)$, where $\pi(x^\pm)=\pi(x)$, so every
$r\in P^\chi_p(w)$ has $\sgn(r)\oplus\pi(r)=\sgn(p)\oplus\pi(p)$. Hence no set
$P^\chi_{p,i}(w)$ meets a flipped set $\overline{P^\chi_{p,j}(w)}$, $\chi$ has no
self events, deadness in a satisfiable $\chi$ is $\DeadT_\chi$, and the live
words of each signed variable form a regular language. This is the automaton
form of \cref{rem:ch-parity}; it includes covariant systems, with $\pi=0$.
\end{remark}

\subsection{Closure Membership}\label{sec:sm-membership}

\begin{lemma}[Closure membership]\label{lem:sm-membership}
Let $(p,w)$ be live in $\chi$, and let $(q,v)$ be an atom. Then
$(q,v)\in\Cl_\chi(p,w)$ iff $(q,v)\in\mu_\chi$ or there are cuts $i\le\len w$
and $e\le\len v$ with
\[
 \wsuf wi\pre\wsuf ve\qquad\text{and}\qquad
 P^\chi_{p,i}(w)\cap Q^\chi_{q,e}(v)\ne\varnothing .
\]
\end{lemma}

\begin{proof}
By \cref{lem:lf-ext-first,lem:sm-lowest-point}, $(q,v)\in \Reach_\chi(p,w)$ iff
there are a word $\alpha$ and cuts $i\le\len{w\alpha}$, $e\le\len v$ with
$\wsuf{(w\alpha)}i=\wsuf ve$ and
$P^\chi_p(\wpre{(w\alpha)}i)\cap Q^\chi_{q,e}(v)\ne\varnothing$. If
$i>\len w$, then $P^\chi_p(w\alpha')\ne\varnothing$ for the nonempty prefix
$\alpha'$ of $\alpha$ of length $i-\len w$, and \cref{lem:sm-roots}(2) with
$k=\len w$ gives $\DeadU_\chi(p,w)$, so $\Cl_\chi(p,w)$ contains a forbidden
root, contradicting liveness by \cref{lem:lf-closure}(2). Hence $i\le\len w$,
and the equation reads $\wsuf wi\,\alpha=\wsuf ve$, which gives
$\wsuf wi\pre\wsuf ve$. Conversely, if $\wsuf wi\,\alpha=\wsuf ve$ and the
intersection is nonempty, \cref{lem:sm-lowest-point} gives
$(p,w\alpha)\to^*_\chi(q,v)$, and $(p,w\alpha)\in \Reach_\chi(p,w)$.
\end{proof}

The two cuts belong to different words, so there is no restriction $i<e$.
The prefix relation points from the seed to the consequence: the seed word may
be continued to the word of the consequence, because extensions are rules.

\subsection{Heads of Derivations}\label{sec:sm-heads}

The pop head of a pure derivation passes only through root-dead atoms.
Immediately before a pop, the signed variable of the current atom is the root
of an upper constructor literal, whose empty-word atom is forbidden.

\begin{lemma}[Head prefixes]\label{lem:sm-head-prefix}
If $P^\chi_p(\gamma)\ne\varnothing$, then $\DeadU_\chi(p,\wpre\gamma i)$ for
every $i<\len\gamma$.
\end{lemma}

\begin{proof}
Since the empty set is absorbing, $P^\chi_p(\wpre\gamma{i+1})\ne\varnothing$,
and \cref{lem:sm-roots}(1) applies to $\wpre\gamma i$ and the letter
$\gamma[i]$.
\end{proof}

The lemma uses neither satisfiability nor any condition relating two systems.
Applied to $\psi$, it says that every proper prefix $\gamma'$ of the pop head
$\gamma$ of a pure derivation of $\psi$ from a visible atom $(p,\gamma\sigma)$
has $\DeadU_\psi(p,\gamma')$; under condition \CondU, the atom $(p,\gamma')$ is
then dead in $\varphi$ (\cref{def:ch-conditions,def:sm-deadu}). \Cref{sec:ex}
uses this to end every pop head of $\psi$ at or before the first rejected cut
of the seed (\cref{lem:ex-fh,lem:ex-head-before-J}).

\subsection{Satisfiability and Condition (I0)}\label{sec:sm-decide}

\begin{proposition}[Satisfiability and (I0)]\label{prop:sm-decide}
\leavevmode
\begin{enumerate}[leftmargin=*]
\item A system $\chi$ is satisfiable iff $\downc\chi m\cap\Fup\chi=\varnothing$
  for every forbidden root $(m,\eps)$ and there are no variable $x$ and word
  $w$ with $(x^+,w),(x^-,w)\in\mu_\chi$.
\item Let $\varphi$ be a system over $X$ and $\psi$ a system over $X\cup Z$.
  Condition \CondI, $\mu_\psi|_X\subseteq\mu_\varphi$, fails iff some
  $q\in X^\pm$ and some word $v$ have $(q,v)\in\mu_\psi$ and
  $(q,v)\notin\mu_\varphi$.
\end{enumerate}
Satisfiability and condition \CondI\ are decidable in polynomial space.
\end{proposition}

\begin{proof}
(1) By \cref{lem:lf-closure}(2), $\chi$ is satisfiable iff $\mu_\chi$ is
consistent, that is, contains no forbidden root and no exclusion pair. By
\cref{lem:sm-mu} with $v=\eps$, $(m,\eps)\in\mu_\chi$ iff
$\downc\chi m\cap\Fup\chi\ne\varnothing$. (2) is the definition of \CondI. For
the complexity, \cref{lem:sm-mu} reads membership in $\mu_\chi$ at $q$ by the
lower automaton of $q$ together with a bit that records whether some state
read so far met $\Fup\chi$. The second condition of (1) is therefore
reachability, for some $x\in V$, of a state with both bits set in the product
of the lower automata of $x^+$ and $x^-$, each with its bit; and the condition
of (2) is reachability, for some $q\in X^\pm$, of a state whose $\psi$-bit is
set and whose $\varphi$-bit is not, in the product of the lower automata of
$q$ in $\psi$ and in $\varphi$. A state of these products has a description of
polynomial size, its successors are computable in polynomial time, and
reachability is decided by guessing a path state by state, in nondeterministic
polynomial space, hence in polynomial space. The set $\Fup\chi$ is computable
in polynomial time by \cref{lem:sm-mu}.
\end{proof}

A shortest witness of the second condition of (1), or of a failure of \CondI,
is shorter than the number of states of the product automaton. Condition
\CondU\ needs more than reachability, because liveness involves self events; it
is the subject of \cref{sec:un}.
